\documentclass[a4paper,14pt]{extarticle}
\usepackage[utf8]{inputenc}
\usepackage[T2A]{fontenc}
\usepackage[russian]{babel}
\usepackage{amsmath}
\usepackage{amsfonts}
\usepackage{amssymb}
\usepackage{listings}
\usepackage{xcolor}
\usepackage{geometry}
\usepackage{hyperref}
\usepackage{booktabs}
\usepackage{tabularx} % Для таблиц с автоматической шириной
\usepackage{array}    % Для дополнительных настроек таблиц

\geometry{left=2cm, right=2cm, top=2cm, bottom=2cm}

% Настройка оформления кода
\definecolor{codegreen}{rgb}{0,0.6,0}
\definecolor{codegray}{rgb}{0.5,0.5,0.5}
\definecolor{codepurple}{rgb}{0.58,0,0.82}
\definecolor{backcolour}{rgb}{0.95,0.95,0.92}
\definecolor{magenta}{rgb}{1,0,1} % Добавлено определение цвета

\lstdefinestyle{mystyle}{
    backgroundcolor=\color{backcolour},
    commentstyle=\color{codegreen},
    keywordstyle=\color{magenta},
    numberstyle=\tiny\color{codegray},
    stringstyle=\color{codepurple},
    basicstyle=\ttfamily\footnotesize,
    breakatwhitespace=false,
    breaklines=true,
    captionpos=b,
    keepspaces=true,
    numbers=left,
    numbersep=5pt,
    showspaces=false,
    showstringspaces=false,
    showtabs=false,
    tabsize=4,
    language=Python
}
\lstset{style=mystyle}

% Настройка для переноса слов в таблицах
\newcolumntype{L}{>{\raggedright\arraybackslash}X}

\begin{document}

\section{Назначение}
\textbf{Сшивка двух изображений} (image stitching) с использованием библиотеки OpenCV.

\section{Структура и основные этапы}

\subsection{Импорт библиотек (ячейка 1)}
\begin{lstlisting}
import cv2
import numpy as np
import matplotlib.pyplot as plt
\end{lstlisting}

\subsection{Функция \texttt{stitch\_images} (ячейка 2)}
Реализует алгоритм сшивки двух изображений:

\begin{table}[h]
\centering
\renewcommand{\arraystretch}{1.3} % Увеличен интервал между строками
\begin{tabularx}{\textwidth}{|l|X|}
\hline
\textbf{Шаг} & \textbf{Описание} \\
\hline
Конвертация в оттенки серого &
\texttt{cv2.cvtColor(..., COLOR\_BGR2GRAY)} \\
\hline
Детекция особенностей & SIFT (\texttt{cv2.SIFT\_create()}) — нахождение ключевых точек и дескрипторов \\
\hline
Сопоставление дескрипторов & \texttt{cv2.BFMatcher()} + \texttt{knnMatch(k=2)} \\
\hline
Фильтрация хороших совпадений &
Тест Лоу (отношение расстояний $< 0.75$) \\
\hline
Вычисление гомографии & \texttt{cv2.findHomography(..., cv2.RANSAC, 5.0)} — оценка матрицы преобразования \\
\hline
Варпирование и сшивка & \texttt{cv2.warpPerspective()} + наложение второго изображения \\
\hline
\end{tabularx}
\caption{Этапы алгоритма сшивки изображений}
\end{table}

\textbf{Возвращает:} сшитое изображение и матрицу гомографии $H$.

\subsection{Загрузка изображений (ячейка 3)}
\begin{lstlisting}
img1 = cv2.imread('k1.jpg')
img2 = cv2.imread('k2.jpg')
\end{lstlisting}

\subsection{Выполнение сшивки (ячейка 4)}
\begin{lstlisting}
result, homography_matrix = stitch_images(img1, img2)
\end{lstlisting}

\subsection{Визуализация результата (ячейка 5)}
Отображение сшитого панорамного изображения через \texttt{matplotlib.pyplot.imshow()}.

\section{Используемые изображения}
\begin{itemize}
\item \texttt{k1.jpg}, \texttt{k2.jpg} — исходные кадры для сшивки (присутствуют в директории проекта)
\end{itemize}

\section{Ключевые алгоритмы и методы}

\begin{table}[h]
\centering
\renewcommand{\arraystretch}{1.3}
\begin{tabularx}{\textwidth}{|l|X|}
\hline
\textbf{Метод} & \textbf{Назначение} \\
\hline
SIFT & Инвариантное к масштабу и повороту детектирование особенностей \\
\hline
BFMatcher + k=2 & Поиск ближайших соседей среди дескрипторов \\
\hline
Lowe's ratio test (0.75) & Отсев ненадёжных совпадений \\
\hline
RANSAC & Робастная оценка матрицы гомографии с отбрасыванием выбросов \\
\hline
warpPerspective & Проективное преобразование первого изображения \\
\hline
\end{tabularx}
\caption{Основные алгоритмы и методы}
\end{table}

\section{Примечания}
\begin{itemize}
\item Код использует \textbf{эвристику 0.75} для фильтрации матчей (стандартное значение из статьи Lowe)
\item Гомография применяется только к первому изображению; второе просто копируется в левую часть результата
\item Возможны артефакты на границе из-за простого наложения без блендинга
\end{itemize}

\end{document}